\documentclass[10pt,a4paper]{amsart}
\usepackage[margin=19mm]{geometry}
\usepackage{amsmath,amssymb,mathtools}
\usepackage[scr=rsfs,cal=euler]{mathalfa}
\usepackage{xfrac}
\usepackage[unicode,hidelinks,bookmarks=false]{hyperref}
\newcommand{\ie}{i.e.\ }
\newcommand{\cf}{cf.\ }
\newcommand{\resp}{respectively\ }
\newcommand{\Subsection}{\subsection}
\providecommand{\nobreakdash}{\nobreak\hbox{-}}
\setlength{\emergencystretch}{2em}
\multlinegap0pt
\usepackage{amssymb}
\usepackage{xcolor}
\usepackage{mathtools}
\usepackage[shortlabels]{enumitem}
\usepackage[normalem]{ulem}
\usepackage{cancel}
\newtheorem{lemma}{Lemma}
\newcommand{\B}{\beta}
\newcommand{\R}{\mathbb{R}}
\newcommand{\blue}{\color{blue}}			
\newcommand{\diff}{d}%{\mathrm{d}}									% Differenziale esterno
\newcommand{\distr}{\mathcal{D}}								% Distribuzione SR
\newcommand{\red}{\color{red}}									
\title{magnetic flows on 3D contact Sub-Riemannian\\ manifolds via the Rumin Complex}
\author{Davide Barilari, Tania Bossio, Valentina Franceschi}
\date{Source: arXiv:2504.09274v2 (CC BY 4.0)}
\begin{document}
\maketitle

\noindent\emph{Source and licence.} magnetic flows on 3D contact Sub-Riemannian\\ manifolds via the Rumin Complex. Version arXiv:2504.09274v2 (CC BY 4.0), distributed by arXiv under the Creative Commons Attribution 4.0 International licence, http://creativecommons.org/licenses/by/4.0/, as recorded in the arXiv licence metadata for this exact version and shown on the abstract page https://arxiv.org/abs/2504.09274v2. Full licence notice: \texttt{licenses/arxiv-2504.09274-CC-BY-4.0.txt}.\par\smallskip

\noindent\textit{Editorial scope.} Counted unit: the article's lemma lem:example (source0.tex lines 1583--1593). The article's complete original proof of it is delivered with it (source0.tex lines 1595--1669, one \texttt{proof} environment of 75 lines). No mathematics is written, repaired or re-derived: the statement and the proof are the article's own lines, in the article's own notation, and no retained line is substituted or re-flowed.  Retained uncounted as the source-local context the proof needs: lines 596--602; lines 674--677; lines 1068--1071.  Nothing was omitted from any retained span, so the package carries no omission record.  No retained line contains a citation, so the wrapper carries no bibliography and no bibliography entry.  No retained line contains a figure environment, an image-inclusion command or drawing code, so no image is delivered and the package declares no asset dependency.  Editorial formatting only, in addition: an AMS \texttt{amsart} preamble that restates the article's own declarations and macros used by the retained lines, namely the environments \texttt{lemma} on separate counters, together with the standard packages those lines need.  The retained environments print in this document as Lemma 1 in order of appearance, whereas the article prints the same items with the numbering of its own full document.  The complete native source of the article as distributed in the arXiv e-print for this exact version is bundled unchanged as \texttt{sources/176429/source0.tex}.
\begin{equation}\label{eq:structurecoeff}
\begin{split}
\left[X_{1},X_{2}\right] & =c_{12}^{1}X_{1}+c_{12}^{2}X_{2}+X_{0}, \\
\left[X_{1},X_{0}\right] & =c_{10}^{1}X_{1}+c_{10}^{2}X_{2},\\
\left[X_{2},X_{0}\right] & =c_{20}^{1}X_{1}+c_{20}^{2}X_{2},
\end{split} 
\end{equation}

\begin{equation}
\label{leibniz}
\mathrm{div}_\mu(fX)=Xf+f\mathrm{div}_\mu(X),
\end{equation}

\begin{equation}
\label{eq:maxwell}
	\mathrm{div}_{\mu}(-\B_2X_1+\B_1X_2)=0.
\end{equation}

\begin{lemma}
\label{lem:example}
Let $n\in\mathbb{N}$, $n\geq 1$ and $f\in C^{\infty}(M)$ be a submersion. Consider a horizontal magnetic field of the form
\begin{equation}
\B=\frac{f^n}{n!}(b_1\nu_1+b_2\nu_2)\wedge\omega\in \Omega^2_H,
\end{equation}
where $b_1,b_2\in C^\infty(M)$ with  $b_1^2+b_2^2 \neq 0.$
%{\red Assume that $\mathrm{rank}\, \diff f=1$}.
Then, the vector field $-b_2X_1+b_1X_2$ is tangent to the zero locus of $\B$ given by $\Sigma=f^{-1}(0)$.
Moreover, $\mathrm{Char}(\Sigma)$ is a 1-dimensional smooth manifold.
\end{lemma}

\begin{proof}
We start by observing that, up to reabsorbing a non zero factor in the function $f$ (thus without changing the zero level set), we can assume that $b_1^2+b_2^2=1$.
Let us consider the orthonormal frame for $\distr$:
\begin{equation}
V_1= b_1X_1+b_2X_2, \qquad V_2 = -b_2X_1+b_1X_2,
\end{equation}
with its corresponding dual basis:
\begin{equation}
\eta_1= b_1\nu_1+b_2\nu_2, \qquad \eta_2=-b_2\nu_1+b_1\nu_2.
\end{equation}
In this basis, we have
\begin{equation}
\B=\frac{f^n}{n!}\eta_1\wedge\omega.
\end{equation}
{\bf Claim (i)}. For every $p\in \Sigma$ we have $V_2f(p)=0$, i.e., the vector field $V_{2}$ is tangent to $\Sigma=f^{-1}(0)$.

\smallskip
Let $\mu$ be the measure associated with the Popp's volume.
Recalling the closure condition \eqref{eq:maxwell}, and applying the Leibniz's rule for the divergence operator in \eqref{leibniz}, we obtain that  
\begin{equation}
0=\mathrm{div}_\mu\left(\frac{f^n}{n!}V_2\right)=\frac{f^{n-1}}{(n-1)!}\left(V_2f+\mathrm{div}_\mu(V_2)\frac{f}{n}\right).
\end{equation}
Therefore, on $M\setminus \Sigma$ (i.e., the set where $f\neq 0$) we have
\begin{equation}
\label{eq:closure}
V_2f+\mathrm{div}_\mu(V_2)\frac{f}{n}=0.
\end{equation} 
Since $\Sigma=f^{-1}(0)$ is a closed set with empty interior, identity \eqref{eq:closure} holds by continuity also on $\Sigma$ and since on $
\Sigma$ we have $f=0$, the Claim (i) is true.

\smallskip
To prove that $\mathrm{Char}(\Sigma)$ is a 1-dimensional smooth manifold, first notice that $\mathrm{Char}(\Sigma)\subset \Sigma$ is the zero level-set of the map $\Phi:M\to \R^2$ defined as $\Phi=(f,V_1f)$. 
Indeed, we have that $p\in \Sigma$ is characteristic if and only if  $V_1f(p)=V_2f(p)=0$ but  $p\in \Sigma$ implies that $V_2f(p)=0$ by \eqref{eq:closure}.

In order to conclude, we show that  for every $p\in \Phi^{-1}(0)$ it holds $\mathrm{rank}\, \diff _p\Phi=2$.
Let us consider the matrix representing $\diff \Phi$ with respect to the frame $V_1,V_2, X_0$ for $TM$:
\begin{equation}\label{eq:db130}
\diff \Phi =
\left( \begin{array}{ccc}
V_1f & V_2f &X_0f \\
V_1V_1f &V_2 V_1f & X_0 V_1f
\end{array} \right).
\end{equation}
{\bf Claim (ii):} Let $p\in \Phi^{-1}(0)$. We have that $V_2V_1f(p)= -X_0f(p)$.

\smallskip
To prove this claim, we first observe that,
differentiating \eqref{eq:closure}, one has
\begin{equation}
\label{eq:V_1V_2f}
V_1 V_2 f =  V_1 \left(-\frac{f}{n}\mathrm{div}_{\mu}(V_2)\right)=-\frac{V_2f}{n}\mathrm{div}_{\mu}(V_2)-\frac{f}{n}V_2\left(\mathrm{div}_{\mu}(V_2)\right).
\end{equation}
For $p\in  \Phi^{-1}(0)$, where $f=V_{2}f=0$, we deduce that $V_1V_2f(p)=0$.

On the other hand, being $V_1,V_2$ an orthonormal frame for $\distr$ with $\diff\omega(V_1,V_2)=-1$, by \eqref{eq:structurecoeff} there exist $a^1_{12},a^2_{12}\in C^\infty(M)$ such that 
\begin{equation}
[V_1,V_2]=a^1_{12}V_1+a^2_{12}V_2+X_0.
\end{equation}
Hence, we obtain that $[V_1,V_2]f(p)=X_0f(p)$ for all $p\in   \Phi^{-1}(0)$. 
Being $V_1V_2f(p)=0$, we deduce
\begin{equation}
\label{eq:commutatorscharS}
V_2V_1f(p)= -X_0f(p),%{\, \blue-\left(b_1^2+b_2^2\right)X_0f(p)}.
\end{equation} 
which proves Claim (ii). Replacing \eqref{eq:commutatorscharS} and $V_1V_2f(p)=0$ in  \eqref{eq:db130}, we have that
\begin{equation}
\diff_p \Phi =
\left( \begin{array}{ccc}
0 & 0 &X_0f(p) \\
V_1V_1f(p) &   -X_0f(p)%{\, \blue-\left(b_1^2+b_2^2\right)X_0f(p)} 
& X_0 V_1f(p)
\end{array} \right).
\end{equation}
Since, by assumption, $f$ is a submersion and $V_1f(p)=V_2f(p)=0$, then necessarily $X_0f(p)\neq 0$. Hence we conclude that $\mathrm{rank}\, \diff _p \Phi =2$, and this concludes the proof.
\end{proof}

\end{document}
